\documentclass[10pt,a4paper]{article}
\usepackage[margin=2.0cm]{geometry}
\usepackage[T1]{fontenc}
\usepackage{lmodern}
\usepackage{microtype}
\usepackage{booktabs}
\usepackage{longtable}
\usepackage{array}
\usepackage{amsmath,amssymb}
\usepackage{xcolor}
\usepackage[hidelinks]{hyperref}
\usepackage{enumitem}
\definecolor{accent}{HTML}{1D4ED8}
\hypersetup{
  colorlinks=true,
  linkcolor=accent,
  urlcolor=accent,
  citecolor=accent,
  pdftitle={Neural MS2LDA: One-Pass Mass2Motif Discovery and Real-Data Model Selection}
}
\newcommand{\StudySeed}{42}
\newcommand{\StudyTopics}{1000}
\newcommand{\StudyThreads}{6}
\newcommand{\FinalEpoch}{40}
\newcommand{\SourceSpectra}{41568}
\newcommand{\RetainedSpectra}{38888}
\newcommand{\TrainingSpectra}{27222}
\newcommand{\ValidationSpectra}{3889}
\newcommand{\TestSpectra}{7777}
\newcommand{\VocabularySize}{21233}
\newcommand{\TomotopyTrainingIterations}{750}
\newcommand{\NeuralValidationOptimized}{843}
\newcommand{\TomotopyValidationOptimized}{607}
\newcommand{\NeuralValidationEvaluable}{429}
\newcommand{\TomotopyValidationEvaluable}{206}
\newcommand{\NeuralValidationUseful}{268}
\newcommand{\TomotopyValidationUseful}{138}
\newcommand{\NeuralValidationSOS}{0.6507}
\newcommand{\TomotopyValidationSOS}{0.6761}
\newcommand{\NeuralValidationMedianSOS}{0.6441}
\newcommand{\TomotopyValidationMedianSOS}{0.6854}
\newcommand{\NeuralTestOptimized}{843}
\newcommand{\TomotopyTestOptimized}{607}
\newcommand{\NeuralTestEvaluable}{567}
\newcommand{\TomotopyTestEvaluable}{333}
\newcommand{\NeuralTestUseful}{347}
\newcommand{\TomotopyTestUseful}{186}
\newcommand{\NeuralTestSOS}{0.6369}
\newcommand{\TomotopyTestSOS}{0.6369}
\newcommand{\NeuralTestMedianSOS}{0.6333}
\newcommand{\TomotopyTestMedianSOS}{0.6145}
\newcommand{\NeuralCoverage}{84.3\%}
\newcommand{\TomotopyCoverage}{60.7\%}
\newcommand{\CoveragePointDelta}{23.6}
\newcommand{\ValidationOptimizedDelta}{236}
\newcommand{\ValidationEvaluableDelta}{223}
\newcommand{\ValidationUsefulDelta}{130}
\newcommand{\TestEvaluableDelta}{234}
\newcommand{\TestUsefulDelta}{161}
\newcommand{\NeuralValidationUsefulShare}{62.5\%}
\newcommand{\TomotopyValidationUsefulShare}{67.0\%}
\newcommand{\NeuralTestUsefulShare}{61.2\%}
\newcommand{\TomotopyTestUsefulShare}{55.9\%}
\newcommand{\NeuralFitMinutes}{35.8}
\newcommand{\TomotopyFitMinutes}{76.0}
\newcommand{\CoreFitRatio}{2.12}
\newcommand{\NeuralActiveTopics}{344}
\newcommand{\NeuralMedianEffectiveTopics}{6.43}
\newcommand{\NeuralValidationNLL}{8.8320}
\newcommand{\TomotopyValidationNLL}{9.6622}
\newcommand{\ValidationNLLDelta}{0.8302}
\newcommand{\NeuralTestNLL}{8.8412}
\newcommand{\TomotopyTestNLL}{9.7569}
\newcommand{\TestNLLDelta}{0.9158}
\newcommand{\WarmBatchSize}{128}
\newcommand{\NeuralWarmRate}{9093.9}
\newcommand{\TomotopyWarmRate}{30.0}
\newcommand{\WarmSpeedRatio}{302.6}

\setlist{nosep,leftmargin=*}
\newcommand{\softmax}{\operatorname{softmax}}
\newcommand{\normalize}{\operatorname{normalize}}
\newcommand{\stopgrad}{\operatorname{stopgrad}}

\title{Neural MS2LDA:\\
One-Pass Mass2Motif Discovery and Real-Data Model Selection}
\author{}
\date{28 August 2026}

\begin{document}
\maketitle

\begin{abstract}
Mass2Motif discovery requires more than good held-out likelihood: a model must
retain a broad, non-redundant topic inventory and assign short, sparse tandem
mass spectra strongly enough for chemical evaluation. Neural MS2LDA (M1) uses
fixed train-only spectral-word embeddings, a contextual token-level top-2
router, shared topic prototypes, a document gate, fragment/loss-balanced
decoding, and complementary anti-collapse objectives. The fitted model has
167,168 learned parameters and performs held-out document inference in one
forward pass.

Under the locked seed-\StudySeed{}, $K=\StudyTopics{}$ MS$^{n}$Lib benchmark,
M1 produces \NeuralValidationOptimized{} optimized and
\NeuralValidationUseful{} useful validation motifs, with completion NLL
\NeuralValidationNLL{}. Against Tomotopy LDA it offers a larger useful motif
inventory and faster resident inference, with a breadth--purity trade-off in
SOS. We then performed a validation-only model-selection campaign against
canonical fixed-SGNS ETM, a deterministic pooled embedded model, a
fragment/loss-balanced ETM, and canonical ECRTM. All simpler completed models
had good completion likelihood, but none preserved M1-level chemistry and topic
inventory. The pooled model contained a 614-topic near-exact duplicate
component; balanced ETM repaired channel skew but not chemical breadth; and the
maintained ECRTM Sinkhorn path failed its convergence contract at real scale.

The evidence supports a qualified conclusion: M1 is not the mathematically
simplest architecture, but it is the least-complex model currently demonstrated
to satisfy the complete real-data scientific contract. Architecture selection
is therefore closed pending multi-seed M1 stability analysis. Alternative model
test data remained locked throughout model selection.
\end{abstract}

\section{Scientific objective}

MS2LDA represents tandem mass spectra as documents, fragment and neutral-loss
tokens as words, and recurrent co-fragmentation patterns as Mass2Motifs
\cite{vanderhooft2016}. At $K=1000$, a useful topic model must satisfy four
requirements simultaneously:
\begin{enumerate}
  \item predict held-out spectral evidence;
  \item retain many distinct topic--word distributions;
  \item assign spectra confidently enough for motif-level association;
  \item yield chemically useful MAG/SOS annotations under leakage control.
\end{enumerate}

These requirements create failure modes that ordinary likelihood can hide.
Topics may be starved, duplicated, or nearly uniform; short spectra may spread
probability over too many topics; and fragment/loss vocabulary geometry may
bias topic mass. A model that predicts documents well can therefore fail as a
Mass2Motif discovery system.

The original M1 study established a neural alternative to Tomotopy and used
within-family ablations to simplify its contextual router. A subsequent concern
was whether M1's combination of routing, gating, Sinkhorn, NPMI, separation, and
alternating optimization was harder to justify than a recognizable published
neural topic model. The validation-only campaign reported here directly tests
that concern.

\section{Neural MS2LDA model}

\subsection{Spectral-word geometry and decoder}

Each training-vocabulary word $w$ has a fixed 50-dimensional feature vector
$x_w$: 48 train-only SGNS coordinates plus fragment/loss indicators. A learned
bias-free projection produces
\begin{equation}
  z_w=\normalize(Wx_w),\qquad W\in\mathbb R^{128\times50}.
\end{equation}
For normalized topic prototype $\widehat q_k$, cosine logits define separate
fragment and loss softmaxes. Every topic assigns one half of its probability
mass to each channel:
\begin{equation}
  \beta_{kw}=\begin{cases}
  \tfrac12\,\softmax_{v\in V_F}(2\widehat q_k^\top z_v/\tau_\beta), & w\in V_F,\\
  \tfrac12\,\softmax_{v\in V_L}(2\widehat q_k^\top z_v/\tau_\beta), & w\in V_L.
  \end{cases}
\end{equation}
The trained topic--word matrix remains an explicit probabilistic Mass2Motif
representation.

\subsection{Contextual sparse routing}

For document $d$, let $c_{dw}$ be the pseudo-count of word $w$ and
$s_d=\sum_wc_{dw}z_w$. The leave-one-out token context is
\begin{equation}
  \bar z_{d\setminus w}=\frac{s_d-c_{dw}z_w}
  {\max(\sum_vc_{dv}-c_{dw},1)}.
\end{equation}
A single bias-free nonlinear residual map combines token and context:
\begin{align}
  e_{dw}&=\operatorname{GELU}(W_c[z_w;\bar z_{d\setminus w}]),\\
  h_{dw}&=\normalize(z_w+e_{dw}),\qquad u_d=\normalize(s_d).
\end{align}
Local and whole-spectrum scores are added, followed by top-2 routing:
\begin{equation}
  a_{dwk}=h_{dw}^{\top}\widehat q_k+u_d^{\top}\widehat q_k,
  \qquad r_{dw}=\operatorname{top2}\!\left(\softmax(a_{dw}/T)\right).
\end{equation}
Count-weighted routed mass $M_{dk}=\sum_wc_{dw}r_{dwk}$ is reweighted by a
detached document gate:
\begin{equation}
  g_{dk}=\softmax_k(u_d^\top\widehat q_k/T),\qquad
  \theta_{dk}=\frac{M_{dk}\,\stopgrad(g_{dk})^{0.75}}
  {\sum_jM_{dj}\,\stopgrad(g_{dj})^{0.75}}.
\end{equation}
The resulting completion model is the ordinary topic mixture
\begin{equation}
  p(w\mid d)=\sum_k\theta_{dk}\beta_{kw}.
\end{equation}
Held-out inference uses this one forward path and no per-spectrum optimization.

\subsection{Training-only safeguards}

M1 alternates router and topic updates under full-spectrum mixture likelihood.
Detached Sinkhorn targets resist topic starvation, positive-NPMI regularization
encourages training-supported co-occurrence, and nearest-prototype separation
resists duplicate topics. Routing temperature and Sinkhorn weight are annealed.
These mechanisms are absent from held-out inference but shape the learned topic
inventory.

The accepted M1 model has 167,168 learned parameters: 6,400 in the token
projection, 32,768 in the context router, and 128,000 in the prototypes. The
complete original derivation and internal simplification campaign are preserved
in
\path{docs/research/archive/neural_ms2lda_report_pre_model_selection.tex}.

\section{Data and evaluation contract}

The experiment uses the published positive-mode MS$^{n}$Lib benchmark assets
from Zenodo record 20179680 \cite{brungs2025,msnlibdata}. The locked seed-42
scaffold/compound-disjoint split contains train, validation, and test partitions;
vocabulary and SGNS features are constructed from training spectra only.
Physical peak groups are divided as units for document completion, preventing a
fragment and its corresponding neutral loss from leaking across observed and
withheld halves.

MAG retrieves leakage-filtered reference spectra and constructs motif
annotations. SOS is then evaluated only for held-out molecules associated with
a topic at $\theta\geq0.5$. The primary chemistry quantities are:
\begin{itemize}
  \item optimized motifs: topics for which MAG yields an optimized motif;
  \item evaluable motifs: optimized topics with at least one high-confidence
  associated held-out molecule;
  \item useful motifs: evaluable topics with SOS $\geq0.6$.
\end{itemize}
Completion NLL is count-weighted over in-vocabulary withheld words. Chemistry
breadth takes precedence over NLL in model selection.

\subsection{Expanded collapse diagnostic contract}

The external comparison showed that optimized motif count alone can be
misleading. Every neural topic-model run must now also report:
\begin{itemize}
  \item median/mean effective topics per spectrum and corpus effective topics;
  \item active-topic counts and maximum mean topic usage;
  \item unique top-1 topics and topics never top-1;
  \item nearest and maximum pairwise beta cosine;
  \item connected duplicate components at cosine 0.95, 0.99, and 0.999;
  \item beta effective-word count, maximum word probability, top-20 mass, and
  top-word uniqueness;
  \item per-topic fragment probability mass and extreme channel skew.
\end{itemize}
These metrics are implemented in
\path{benchmarks/neural_ms2lda/diagnostics.py} and are written for validation and
test evaluation. Candidate selection remains validation-only.

\section{Locked M1 versus Tomotopy result}

On validation, M1 produces \NeuralValidationOptimized{} optimized,
\NeuralValidationEvaluable{} evaluable, and \NeuralValidationUseful{} useful
motifs. Tomotopy produces \TomotopyValidationOptimized{},
\TomotopyValidationEvaluable{}, and \TomotopyValidationUseful{}, respectively.
M1 has lower validation completion NLL
(\NeuralValidationNLL{} versus \TomotopyValidationNLL{}), while Tomotopy has
higher validation mean SOS
(\TomotopyValidationSOS{} versus \NeuralValidationSOS{}). Thus the validation
result is a breadth--purity trade-off rather than uniform dominance.

After model selection and lock, the original study opened test once for
reporting. M1 yields \NeuralTestUseful{} useful test motifs versus
\TomotopyTestUseful{} for Tomotopy, with completion NLL \NeuralTestNLL{} versus
\TomotopyTestNLL{}. The recorded core fitting times are \NeuralFitMinutes{} and
\TomotopyFitMinutes{} minutes. These test numbers apply only to the already
locked M1/Tomotopy comparison; no later candidate was evaluated on test.

\section{Validation-only external model selection}

The comparison started from recognizable published formulations and added only
failure-motivated adaptations. The frozen candidate gates were optimized
$\geq840$, evaluable $\geq388$, useful $\geq252$, mean SOS $\geq0.651498$,
completion NLL $\leq9.422847$, finite execution, and no catastrophic inventory
collapse.

\begin{table}[ht]
\centering
\scriptsize
\setlength{\tabcolsep}{3.5pt}
\caption{Real MS$^{n}$Lib validation-only model-selection result. Alternatives
were never evaluated on candidate test data.}
\label{tab:model-selection}
\begin{tabular}{@{}>{\raggedright\arraybackslash}p{0.25\textwidth}rrrrr>{\raggedright\arraybackslash}p{0.19\textwidth}@{}}
\toprule
Model & Optimized & Evaluable & Useful & Mean SOS & NLL & Outcome \\
\midrule
M1 locked reference & 884 & 408 & 265 & 0.658079 & 8.974140 & Pass \\
Canonical fixed-SGNS ETM & 609 & 130 & 79 & 0.638796 & 8.690730 & Fail chemistry \\
Pooled projected & 967 & 14 & 11 & 0.751621 & 8.385787 & Fail breadth/collapse \\
Pooled projected, $\tau=0.11$ & 967 & 293 & 194 & 0.665436 & 8.994812 & Fail breadth/collapse \\
Pooled projected + MI 0.05 & 799 & 14 & 11 & 0.696540 & 8.393690 & Fail breadth/collapse \\
Fragment/loss-balanced ETM & 911 & 166 & 104 & 0.629819 & 8.766069 & Fail breadth/SOS \\
Canonical ECRTM & -- & -- & -- & -- & -- & Infeasible at epoch 22 \\
\bottomrule

\end{tabular}
\end{table}

\subsection{Canonical fixed-SGNS ETM}

The faithful ETM uses a two-layer variational document encoder and a global
embedding-space topic--word softmax \cite{dieng2020}. Fixed spectral SGNS avoided
the severe collapse observed with jointly learned embeddings in simulation.
On real validation, ETM achieved lower NLL than M1 but only 609 optimized and 79
useful motifs. Its median spectrum used 43.79 effective topics. Temperature
sharpening increased associations but could not repair the topic-word deficit.

\subsection{Pooled projected model}

The deterministic pooled model projects token features, count-pools one spectrum
vector, and uses a shared prototype bank for both $\theta$ and $\beta$. Its raw
NLL and MAG optimization coverage appeared excellent, but its document mixtures
were too diffuse. Post-hoc temperature $\tau=0.11$ matched M1-like mixture
sparsity and raised evaluable/useful counts from 14/11 to 293/194.

The deeper diagnosis showed that calibration was not the whole problem. A
single connected component contained 614 topics at beta cosine $\geq0.999$.
Those topics were nearly uniform over the vocabulary, never won a validation
spectrum, and could not be separated by rank-preserving temperature changes.
Weak assignment MI left essentially the same collapse.

\subsection{Fragment/loss-balanced ETM}

Real canonical ETM learned substantial channel skew, so a paired ETM was trained
with independent 50/50 fragment/loss softmaxes. This adaptation eliminated the
skew and raised optimized motifs from 609 to 911, demonstrating that global
normalization was a real source of topic-spectrum distortion. Nevertheless,
evaluable/useful motifs remained 166/104 and mean SOS decreased. Channel
balance repaired annotation coverage but not the complete Mass2Motif objective.

\subsection{ECRTM}

ECRTM is a published embedding-clustering regularized topic model designed for
component collapse \cite{wu2023}. It was scientifically justified after the
614-topic pooled component was observed. At real $K=1000$ and vocabulary
21,233, the maintained ordinary-domain Sinkhorn solver became progressively
more demanding. After 21 completed epochs, mean/max iterations reached 721/951;
epoch 22 failed the residual tolerance within 1,000 iterations and an exact
checkpoint resume reproduced the failure. No partial or unconverged model was
scored. This is an operational conclusion about the maintained numerical path,
not a universal impossibility result for stabilized ECRTM implementations.

\begin{table}[ht]
\centering
\scriptsize
\caption{Diagnostics explaining why good likelihood did not imply a usable
Mass2Motif inventory. The largest component uses beta cosine $\geq0.999$.}
\label{tab:model-diagnostics}
\begin{tabular}{@{}>{\raggedright\arraybackslash}p{0.29\textwidth}rrrrr@{}}
\toprule
Model & Median effective topics & Unique top-1 & Largest component & Mean nearest cosine & Max cosine \\
\midrule
Canonical fixed-SGNS ETM & 43.79 & 260 & 2 & 0.321 & 0.9995 \\
Pooled projected & 130.01 & 374 & 614 & 0.692 & 1.0000 \\
Pooled projected, $\tau=0.11$ & 7.35 & 374 & 614 & 0.692 & 1.0000 \\
Fragment/loss-balanced ETM & 46.72 & 260 & 1 & 0.400 & 0.9958 \\
\bottomrule

\end{tabular}
\end{table}

\section{Interpretation}

\subsection{Why M1's complexity is functional}

The external failures align with the roles of M1's mechanisms. The pooled model
demonstrates severe prototype/topic duplication; ETM demonstrates diffuse and
weak spectrum-topic assignment; balanced ETM shows that annotatable topic
spectra do not by themselves guarantee chemical specificity. M1's contextual
sparse routing and document gate address assignment, Sinkhorn addresses usage,
prototype separation addresses duplicate components, positive NPMI supports
coherent words, and the channel split controls fragment/loss mass.

Within-family M1 ablations had already shown substantial motif losses when the
document gate, Sinkhorn, NPMI, or prototype separation was removed. The external
campaign now shows that moving to a simpler model family does not make those
failure modes disappear. M1 should therefore be described as the least-complex
model demonstrated to meet the full contract, not as the simplest conceivable
architecture or as a uniformly superior generic topic model.

\subsection{What is not claimed}

The model-selection campaign uses one positive-mode MS$^{n}$Lib benchmark and
one locked data split. It does not prove that M1 is optimal on other instruments,
ionization modes, biological samples, or topic counts. It also does not rule out
a future stabilized ECRTM, a different published diversity regularizer, or a
new representation. It establishes that the completed alternatives tested here
do not justify replacing M1.

MAG/SOS remains a library-dependent chemical proxy. A high optimized count can
include broad or unused topics, which is why the expanded diagnostic contract is
now mandatory. Expert interpretation and independent-dataset replication remain
necessary.

\section{Next evidence: M1 stability}

Architecture search is paused. The next compute campaign varies only M1
optimization seed while keeping the seed-42 split, vocabulary, SGNS features,
co-occurrence graph, and validation contract byte-identical. New training seeds
11, 23, and 37 are predeclared; test remains locked. The campaign will report
per-seed chemistry, completion, the expanded collapse diagnostics, and
one-to-one cross-seed beta alignment. The exact handoff is in
\path{research/etm_ecrtm_msnlib/M1_MULTISEED_HANDOFF.md}.

This design separates optimization stability from data-split variability. A
later independent-library study should examine broader generalization, but it
should not be conflated with the immediate question of whether the seed-42 M1
result is representative of repeated fitting on the same benchmark.

\section{Conclusion}

Neural MS2LDA performs one-pass document inference while retaining explicit
probabilistic Mass2Motifs. On the locked MS$^{n}$Lib benchmark it yields a broad,
chemically useful inventory and strong completion performance relative to
Tomotopy. A subsequent validation-only comparison showed that canonical ETM,
a deterministic pooled embedded model, balanced ETM, and canonical ECRTM each
failed a different part of the scientific contract.

The combined evidence changes the justification for M1. Its architecture is not
presented as novelty for novelty's sake; its mechanisms correspond to observed
real-data failure modes that simpler models did not jointly solve. M1 remains
the only gate-passing neural model and no new candidate advances to test. The
next claim to establish is reproducibility across optimization seeds, not another
unbounded architecture search.

\appendix
\section{Reproduction and provenance}

Generate the original numerical macros and the model-selection tables with:
\begin{verbatim}
python scripts/generate_neural_ms2lda_report.py
python scripts/generate_neural_ms2lda_model_selection.py
\end{verbatim}
Then compile \path{docs/research/neural_ms2lda_report.tex}. The model-selection
source is the committed validation-only comparison at
\path{research/etm_ecrtm_msnlib/local_results/20260827_followup/comparison.csv}.
The full hypothesis-by-hypothesis record is in the neighboring
\path{EXPERIMENT_LOG.md}. The previous detailed methods report is archived
verbatim rather than deleted.

\section{Frozen diagnostic settings}

The reporting contract uses mean topic usage 0.0005, duplicate-component cosine
thresholds 0.95/0.99/0.999, top 20 words, catastrophic strict-component fraction
0.5, and fragment-mass extremes below 0.1 or above 0.9. These are reporting and
failure-detection settings; chemistry selection still uses the separately
frozen MAG/SOS gates.

\begin{thebibliography}{99}
\bibitem{vanderhooft2016}
J. J. J. van der Hooft et al. ``Topic modeling for untargeted substructure
exploration in metabolomics.'' \emph{Proceedings of the National Academy of
Sciences} 113, 13738--13743 (2016).
\url{https://doi.org/10.1073/pnas.1608041113}

\bibitem{torresortega2026}
M. Torres Ortega et al. ``Large-scale discovery and annotation of substructure
patterns in mass spectrometry profiles.'' \emph{Nature Communications} 17,
8350 (2026). \url{https://doi.org/10.1038/s41467-026-75038-0}

\bibitem{brungs2025}
C. Brungs et al. ``MS$^{n}$Lib: efficient generation of open multi-stage
fragmentation mass spectral libraries.'' \emph{Nature Methods} 22, 2028--2031
(2025). \url{https://doi.org/10.1038/s41592-025-02813-0}

\bibitem{dieng2020}
A. B. Dieng, F. J. R. Ruiz, and D. M. Blei. ``Topic Modeling in Embedding
Spaces.'' \emph{Transactions of the Association for Computational Linguistics}
8, 439--453 (2020). \url{https://doi.org/10.1162/tacl_a_00325}

\bibitem{wu2023}
X. Wu, X. Dong, T. Nguyen, and A. T. Luu. ``Effective Neural Topic Modeling
with Embedding Clustering Regularization.'' \emph{Proceedings of the 40th
International Conference on Machine Learning}, 2023.
\url{https://proceedings.mlr.press/v202/wu23c.html}

\bibitem{msnlibdata}
J. Dietrich, J. Wandy, L. R. Torres Ortega, and J. J. J. van der Hooft.
\emph{vdhooftcompmet/MS2LDA: Paper---Data and Models}. Zenodo (2026).
\url{https://doi.org/10.5281/zenodo.20179680}
\end{thebibliography}

\end{document}
